\documentclass[10pt,a4paper]{amsart}
\usepackage[margin=19mm]{geometry}
\usepackage{amsmath,amssymb,mathtools}
\usepackage[scr=rsfs,cal=euler]{mathalfa}
\usepackage{xfrac}
\usepackage[unicode,hidelinks,bookmarks=false]{hyperref}
\newcommand{\ie}{i.e.\ }
\newcommand{\cf}{cf.\ }
\newcommand{\resp}{respectively\ }
\newcommand{\Subsection}{\subsection}
\providecommand{\nobreakdash}{\nobreak\hbox{-}}
\setlength{\emergencystretch}{2em}
\multlinegap0pt
\usepackage{bm}
\usepackage{bbm}
\usepackage{dsfont}
\usepackage{xspace}
\usepackage{cancel}
\usepackage{mathtools}
\usepackage{amsthm}
\makeatletter
\@ifundefined{theorem}{\newtheorem{theorem}{Theorem}}{}
\makeatother
\title[Multiplicity for partially ordered sets]{Multiplicity for partially ordered sets}
\author{Gyula O.H. Katona, Yaping Mao}
\date{Source: arXiv:2607.00456v1 (CC BY 4.0)}
\begin{document}
\maketitle

\noindent\emph{Source and licence.} Multiplicity for partially ordered sets. Version arXiv:2607.00456v1 (CC BY 4.0), distributed under the Creative Commons Attribution 4.0 International licence, as recorded in the arXiv licence metadata for this exact version and shown on the abstract page https://arxiv.org/abs/2607.00456v1. Full rights notice: licenses/arxiv-2607.00456-CC-BY-4.0.txt.\par\smallskip

\noindent\textit{Editorial scope.} Counted unit: the Theorem whose source label is \texttt{thm:main} in the article (source0.tex lines 215--247, source label \texttt{thm:main}), delivered together with the article's own complete original proof of it (source0.tex lines 320--450, one \texttt{proof} environment of 131 lines). No mathematics is written, repaired or re-derived: statement and proof are the article's own text line for line. The proof is detached from the statement in the source: the article states the result at lines 215--247 and prints its proof at lines 320--450, and the proof environment's own header names the statement, so the association is the article's own. Required context retained uncounted: none. Every definition, display and label of those lines is retained unchanged. Omitted from this package and not redistributed: 1--214, 248--319, 451--1334. Nothing inside a retained excerpt is omitted or altered: every retained body is delivered exactly as the article's lines are written, so no omission receipt of any kind accompanies this package. Citations: the retained excerpts carry no citation command at all, so no bibliography entry of the article is transcribed and no reference is redistributed. Reference adaptation: none was needed and none was made; every reference inside the delivered unit resolves inside it, so no unresolved reference or citation is produced. Printed slips kept literal: no printed word, symbol, hypothesis or number of the counted statement or of its proof was altered. Layout and class: the article's own class is \texttt{article} with packages including inputenc, mathrsfs, amssymb, amsmath, amsthm, graphicx, color, tikz, geometry, float, caption, subcaption, diagbox, algorithm; the wrapper is the lane's pinned \texttt{amsart} editorial wrapper at 10pt on A4, which restates the theorem-like environments the retained text needs and adds the article's own macro definitions that the retained text needs (none), with extra wrapper packages (bm, bbm, dsfont, xspace, cancel, mathtools). The delivered numbering is the unit's own. Assets: no retained span contains a figure environment, a graphics-inclusion command, a PDF-page inclusion, a table or a TikZ picture; no image file is copied into this package, the package declares no asset dependency and no image file is redistributed with it. Licence: version arXiv:2607.00456v1 of this article is distributed under the Creative Commons Attribution 4.0 International licence, as recorded in the arXiv licence metadata for this exact version and shown on the abstract page https://arxiv.org/abs/2607.00456v1. A full rights notice is delivered at licenses/arxiv-2607.00456-CC-BY-4.0.txt. Characters: every retained excerpt is pure ASCII, so the delivered engine, XeLaTeX with its default Latin Modern text font, reports no missing character.
\begin{theorem}\label{thm:main}
Let
\[
 E_n=\{(S,T,U)\in B_n^3:S\subsetneq T\subsetneq U,\ |S|+|T|=|U|\}.
\]
Then the two-color threshold for forcing a monochromatic member of $E_n$ equals $9$. Moreover,
\[
|E_n|
=\binom{2n}{n}-[x^n](1+x+x^2)^n-2^n+1
=\frac{4^n}{\sqrt{\pi n}}\bigl(1+o(1)\bigr).
\]
If $M^{\mathrm{arith}}_2(B_n)$ denotes the minimum number of monochromatic members of $E_n$ over all two-colorings of $B_n$, then
\[
2^{\delta n+o(n)}
\leq M^{\mathrm{arith}}_2(B_n)
\leq 2^{\gamma n+o(n)},
\]
where
\[
\delta=-3\cdot\frac1{10}\log_2\frac1{10}-\frac7{10}\log_2\frac7{10}
      \approx 1.356779
\]
and
\[
\gamma=\min_{1/3\leq\alpha\leq1/2}
\max\{H(\alpha)+\alpha\log_2 3,\ 2H(1-2\alpha)\}
      \approx 1.567837.
\]
Here $H(x)=-x\log_2 x-(1-x)\log_2(1-x)$, with $H(0)=0$ by continuity, and the minimum is attained at the unique solution $\alpha_0\approx0.383292$ of
\[
H(\alpha)+\alpha\log_2 3=2H(1-2\alpha).
\]
\end{theorem}

\begin{proof}[Proof of Theorem \ref{thm:main}]
Fix ranks $|S|=k$, $|T|=m$, and $|U|=k+m$. Since
$S\subsetneq T\subsetneq U$,
we must have $1\leq k<m$, and the condition $|U|=k+m$ gives $k+m\leq n$.

For fixed $k$ and $m$, we count the triples $(S,T,U)$ as follows. First choose the middle set $T$ in $\binom{n}{m}$ ways. Then choose $S\subset T$ with $|S|=k$ in $\binom{m}{k}$ ways. Finally, since $U$ must contain $T$ and have size $k+m$, the set $U\setminus T$ must be a $k$-subset of $[n]\setminus T$, and this can be chosen in $\binom{n-m}{k}$ ways. Hence
\[
|E_n|
=
\sum_{\substack{1\leq k<m\\ k+m\leq n}}
\binom{n}{m}\binom{m}{k}\binom{n-m}{k}.
\]
We now transform this sum. Using
\[
\binom{n}{m}\binom{m}{k}
=
\binom{n}{k}\binom{n-k}{m-k},
\]
we obtain
\begin{align*}
\sum_{m=k+1}^{n-k}\binom{n}{m}\binom{m}{k}\binom{n-m}{k}
&=
\binom{n}{k}
\sum_{m=k+1}^{n-k}
\binom{n-k}{m-k}\binom{n-m}{k}.
\end{align*}
Let $\ell=m-k$. Then $1\leq \ell\leq n-2k$, and therefore
\begin{align*}
\sum_{m=k+1}^{n-k}\binom{n}{m}\binom{m}{k}\binom{n-m}{k}
&=
\binom{n}{k}
\sum_{\ell=1}^{n-2k}
\binom{n-k}{\ell}\binom{n-k-\ell}{k}  \\
&=
\binom{n}{k}\binom{n-k}{k}
\sum_{\ell=1}^{n-2k}
\binom{n-2k}{\ell}  \\
&=
\binom{n}{k}\binom{n-k}{k}
\bigl(2^{n-2k}-1\bigr).
\end{align*}
Summing over $1\leq k\leq \lfloor n/2\rfloor$ gives
\[
|E_n|
=
\sum_{k=1}^{\lfloor n/2\rfloor}
\binom{n}{k}\binom{n-k}{k}
\bigl(2^{n-2k}-1\bigr).
\]
It remains to derive the equivalent form. Let
\[
A_n=
\sum_{k=0}^{\lfloor n/2\rfloor}
\binom{n}{k}\binom{n-k}{k}2^{n-2k}
\]
and
\[
T_n=
\sum_{k=0}^{\lfloor n/2\rfloor}
\binom{n}{k}\binom{n-k}{k}.
\]
Then the preceding formula gives
$|E_n|=A_n-T_n-2^n+1$,
because the $k=0$ term of $A_n$ is $2^n$, while the $k=0$ term of $T_n$ is $1$.

We claim that
$A_n=\binom{2n}{n}$.
Indeed, consider a set of \(2n\) distinct elements divided into \(n\) fixed
pairs. We count the number of ways to choose \(n\) elements from this set.
Clearly, the answer is \(\binom{2n}{n}\).
We count the same quantity according to the number of pairs from which both
elements are chosen. Suppose exactly \(k\) pairs contribute both of their
elements. Since the chosen set has size \(n\), exactly \(k\) pairs must
contribute no element, and the remaining \(n-2k\) pairs contribute exactly
one element each. The \(k\) pairs contributing two elements can be chosen in
\(\binom{n}{k}\) ways. Then the \(k\) pairs contributing no element can be
chosen in \(\binom{n-k}{k}\) ways. Finally, from each of the remaining
\(n-2k\) pairs, one of the two elements is chosen, giving \(2^{n-2k}\)
choices. Hence, for fixed \(k\), the number of choices is
\[
\binom{n}{k}\binom{n-k}{k}2^{n-2k}.
\]
Summing over all possible \(k\), we obtain
\[
A_n
=
\sum_{k=0}^{\lfloor n/2\rfloor}
\binom{n}{k}\binom{n-k}{k}2^{n-2k}
=
\binom{2n}{n}.
\]

Therefore
\[
|E_n|
=
\binom{2n}{n}
-
T_n
-
2^n+1
=
\binom{2n}{n}
-
\sum_{k=0}^{\lfloor n/2\rfloor}
\binom{n}{k}\binom{n-k}{k}
-
2^n+1.
\]

Finally, we estimate the error term. The quantity $T_n$ counts ordered pairs $(A,B)$ of disjoint subsets of $[n]$ with $|A|=|B|$. Hence $T_n$ is at most the number of all ordered pairs of disjoint subsets of $[n]$, which is $3^n$, since each element of $[n]$ may lie in $A$, lie in $B$, or lie in neither. Thus
$T_n\leq 3^n$.
Consequently,
\[
|E_n|
=
\binom{2n}{n}+O(3^n).
\]
Using the standard asymptotic formula
\[
\binom{2n}{n}
=
\frac{4^n}{\sqrt{\pi n}}\bigl(1+o(1)\bigr),
\]
and noting that $3^n=o(4^n/\sqrt n)$, we obtain that 
\[
|E_n|
=
\frac{4^n}{\sqrt{\pi n}}\bigl(1+o(1)\bigr).
\]
\end{proof}

\end{document}
